% !TeX program = xelatex
\documentclass[aspectratio=169, 11pt]{beamer}
\usepackage[UTF8]{ctex}
\usepackage{seu_clean_v2}
\usepackage{booktabs}
\usepackage{graphicx}

\title[简洁版模板示例]{东南大学 Beamer 简洁版模板 \\ Example of Clean Template}
\subtitle{顶边绿字 + 正文白底 + CJK 支持}
\author[演示作者]{演示作者 \\ Southeast University}
\date{\today}

\begin{document}

\begin{frame}
\titlepage
\end{frame}

\section{模板特性概览}
\section{排版元素演示}
\section{表格与公式}
\section{总结}

\plainframetitle
\begin{frame}{目\quad 录}
\tableofcontents
\end{frame}

\begin{frame}{简洁版模板特性}
\begin{block}{设计特点}
  \begin{itemize}
    \item 仅标题页有 SEU 背景水印，正文页纯白背景
    \item 顶边文字为 SEU 绿色（\textcolor{structure.fg}{structure.fg}），与底边统一
    \item 顶边背景浅绿底色，字号 \texttt{\textbackslash footnotesize}
    \item CJK 中文支持（xeCJK + Microsoft YaHei）
    \item 保留标题页布局（顶边不可见占位）
  \end{itemize}
\end{block}

\begin{exampleblock}{适用场景}
  国际会议、学术报告、非正式分享、偏爱清爽视觉风格的场合。
\end{exampleblock}
\end{frame}

\section{排版元素演示}

\begin{frame}{列表与枚举}
\begin{columns}
\column{0.48\textwidth}
\begin{block}{itemize 列表}
  \begin{itemize}
    \item 第一级项目
    \item 第一级项目
    \begin{itemize}
      \item 第二级子项目
      \item 第二级子项目
    \end{itemize}
  \end{itemize}
\end{block}

\column{0.48\textwidth}
\begin{block}{enumerate 列表}
  \begin{enumerate}
    \item 第一步操作
    \item 第二步操作
    \item 第三步操作
  \end{enumerate}
\end{block}
\end{columns}
\end{frame}

\begin{frame}{Block 环境演示}
\begin{block}{普通 Block}
  这是一般性重要信息。用于正文中的关键内容强调。
\end{block}

\vspace{0.3cm}

\begin{alertblock}{警告 Block}
  用于关键警告或注意事项。每帧最多一个。
\end{alertblock}

\vspace{0.3cm}

\begin{exampleblock}{示例 Block}
  用于示例、案例或直觉解释。每帧最多一个。
\end{exampleblock}
\end{frame}

\begin{frame}{双栏布局}
\begin{columns}
\column{0.45\textwidth}
\centering
\textbf{左栏内容}\\
\vspace{0.3cm}
\footnotesize
左栏可以放置文字说明、列表或简要描述。\\
在双栏布局中，左右栏的内容应均衡，避免一栏过密另一栏过空。

\column{0.45\textwidth}
\centering
\textbf{右栏内容}\\
\vspace{0.3cm}
\footnotesize
右栏可以放置补充说明、公式或图片。\\
columns 环境是 Beamer 中实现并排布局的标准方式。
\end{columns}
\end{frame}

\section{表格与公式}

\begin{frame}{表格示例（booktabs）}
\begin{table}
\centering
\caption{各方法性能对比}
\begin{tabular}{lccc}
\toprule
\textbf{方法} & \textbf{准确率} & \textbf{召回率} & \textbf{F1} \\
\midrule
方法 A & 89.2\% & 87.6\% & 88.4 \\
方法 B & 94.1\% & 92.3\% & 93.2 \\
方法 C & \textbf{95.7\%} & \textbf{94.8\%} & \textbf{95.2} \\
\bottomrule
\end{tabular}
\end{table}

\vspace{0.2cm}

\footnotesize 表格使用 \texttt{booktabs} 宏包规范，禁止竖线，表头加粗。
\end{frame}

\begin{frame}{公式与数学}
\begin{block}{核心公式}
  \begin{equation}
    \mathcal{L}(\theta) = -\sum_{i=1}^{n} \log P(y_i \mid x_i; \theta) + \lambda \|\theta\|_2^2
  \end{equation}
\end{block}

\vspace{0.2cm}

\begin{itemize}
  \item 交叉熵损失 + L2 正则化
  \item 支持 \texttt{amsmath} 全部数学环境
\end{itemize}
\end{frame}

\section{总结}

\begin{frame}{模板使用方式}
\begin{enumerate}
  \item 在 \texttt{SEU-Beamer-Slide/} 目录下编写 \texttt{.tex} 文件
  \item 导言区使用 \texttt{\textbackslash usepackage\{seu\textunderscore clean\}}
  \item 使用 \texttt{xelatex} 编译两次
  \item 将生成的 PDF 复制到目标位置
\end{enumerate}

\vspace{0.3cm}

\begin{alertblock}{注意}
  编译前确保 \texttt{source/} 目录（含 SEU logo 图片）与 \texttt{.tex} 文件在同一目录下。
\end{alertblock}
\end{frame}

\end{document}
